% from ega0-20-fr.tex
\oldpage[0\textsubscript{IV-1}]{119}
\begin{env}[20.2.1]
\label{0.20.2.1-fr}
Soient $A$ un anneau, $B$ un $A$-anneau, $L$ un $B$-bimodule ; si $L'$ est
un second $B$-bimodule et $w:L\to L'$ un homomorphisme de $B$-bimodules, il
est clair que l'application $D\mapsto w\circ D$ est un homomorphisme de
groupes additifs
\begin{equation}
\tag{20.2.1.1}
\label{0.20.2.1.1-fr}
  w_0:\operatorname{D\acute{e}r}_A(B,L)\to\operatorname{D\acute{e}r}_A(B,L')
\end{equation}
et que si $w':L'\to L''$ est un second homomorphisme de $B$-bimodules, on a
$(w'\circ w)_0=w'_0\circ w_0$. Lorsque $A$ est commutatif, $B$ une
$A$-algèbre commutative et $L$ un $B$-module, \sref{0.20.2.1.1-fr} est un
homomorphisme de $A$-modules.

En second lieu, soient $B'$ un $A$-anneau, $v:B'\to B$ un
$A$-homomorphisme qui fait de $L$ un $B'$-bimodule ; alors l'application
$D\mapsto D\circ v$ est un homomorphisme de groupes additifs,
\begin{equation}
\tag{20.2.1.2}
\label{0.20.2.1.2-fr}
  v^0:\operatorname{D\acute{e}r}_A(B,L)\to\operatorname{D\acute{e}r}_A(B',L)
\end{equation}
comme il résulte de \sref{0.20.1.1.1-fr} ; si $v':B''\to B'$ est un second
$A$-homomorphisme, on a $(v\circ v')^0=v'^0\circ v^0$. Lorsque $A$, $B$ et
$B'$ sont commutatifs et $L$ un $B$-module, \sref{0.20.2.1.2-fr} est un
homomorphisme de $A$-modules.

Enfin, soit $u:A'\to A$ un homomorphisme d'anneaux faisant de $B$ un
$A'$-anneau ; toute $A$-dérivation $D\in\operatorname{D\acute{e}r}_A(B,L)$ est aussi
une $A'$-dérivation, d'où une injection canonique de groupes commutatifs
\begin{equation}
\tag{20.2.1.3}
\label{0.20.2.1.3-fr}
  u^0:\operatorname{D\acute{e}r}_A(B,L)\to\operatorname{D\acute{e}r}_{A'}(B,L)
\end{equation}
et si $u':A''\to A'$ est un second homomorphisme d'anneaux, on a
$(u\circ u')^0=u'^0\circ u^0$ ; lorsque $A$, $A'$ et $B$ sont commutatifs et
$L$ un $B$-module, \sref{0.20.2.1.3-fr} est un di-homomorphisme de modules
(relatif à $u$).

On peut encore dire que
\[
  (A,B,L)\leadsto\operatorname{D\acute{e}r}_A(B,L)
\]
est un \emph{foncteur covariant} de la catégorie $\boldsymbol{K}$ définie dans
(18.3.5) dans la catégorie $\boldsymbol{Ab}$ des groupes commutatifs, en
faisant correspondre, à tout triplet $(u,v,w)$ constituant un morphisme de
$\boldsymbol{K}$ l'homomorphisme $w_0\circ v^0\circ u^0$ ; la vérification de
la fonctorialité résulte de la commutativité des diagrammes
\[
\xymatrix{
  \operatorname{D\acute{e}r}_A(B,L) \ar[r]^{v^0} \ar[d]_{w_0} &
    \operatorname{D\acute{e}r}_A(B',L) \ar[d]^{w_0} \\
  \operatorname{D\acute{e}r}_A(B,L') \ar[r]^{v^0} &
    \operatorname{D\acute{e}r}_A(B',L')
}
\qquad
\xymatrix{
  \operatorname{D\acute{e}r}_A(B,L) \ar[r]^{u^0} \ar[d]_{w_0} &
    \operatorname{D\acute{e}r}_{A'}(B,L) \ar[d]^{w_0} \\
  \operatorname{D\acute{e}r}_A(B,L') \ar[r]^{u^0} &
    \operatorname{D\acute{e}r}_{A'}(B,L')
}
\]
pour tout homomorphisme $w:L\to L'$ de $B$-bimodules.
\end{env}

